%%%%%%%%%%%%%%%%%%%%%%%%%%%%%%
%%%%%%%%%%%% Preamble %%%%%%%%%%%%%
%%%%%%%%%%%%%%%%%%%%%%%%%%%%%%

% Declare document class and miscellaneous packages
\documentclass[english]{article}
\usepackage{natbib}
\usepackage[leqno]{amsmath}
\usepackage{amssymb}
\usepackage{mathrsfs}
\usepackage{mathtools}
\usepackage{setspace}
\usepackage{fullpage}
\usepackage[margin=1in]{geometry}
\interfootnotelinepenalty=10000

%Hyper-references
\usepackage{hyperref}
\hypersetup{colorlinks, citecolor=black, filecolor=black, linkcolor=black, 
urlcolor=black, pdftex}

\begin{document}

\doublespacing

%%%%%%%%%%%%%%%%%%%%%%%%%%%%%%
%%%%%%%%%%%%% Body %%%%%%%%%%%%%%
%%%%%%%%%%%%%%%%%%%%%%%%%%%%%%

% This write-up is based on the post at the following URL:
% https://stats.stackexchange.com/questions/238878/how-do-errors-in-variables-affect-the-r2

Let $y_i^t$ denote the true poverty rate for some geographical entity $i = 1, 2, \ldots, N$.
Further, let $y_i^d$ and $y_i^m$ denote the direct and model-based estimates of the poverty rate, 
respectively.
The standard method for assessing the predictive performance of model-based poverty estimates is
to calculate the $R^2$ from a regression like the following:
\begin{equation}
y_i^m = \alpha + \beta y_i^d + \varepsilon_i
\label{eq:reg}
\end{equation}
where we emphasize that the above is a regression of the model-based estimates on the direct
estimates rather than the true poverty rates.
The relationship between the direct estimates and the true poverty rates can be expressed as follows:
\begin{equation}
y_i^d = y_i^t + \eta_i
\end{equation}
where we assume that $\eta$ is a mean-zero error term that is independent of $y_i^t$.
We further assume that $\eta$ and $\varepsilon$ are independent.

The estimated $R^2$ from the above regression can be written as
\begin{equation}
\hat{R}^2 = 1 - \frac{\sum_i(\hat{\alpha} + \hat{\beta}y_i^d - y_i^m)^2}{\sum_i (y_i^m - \bar{y}^m)^2}
\end{equation}
where $\hat{\alpha}$ and $\hat{\beta}$ are the coefficient estimates.
Noting that $\hat{\alpha} = \bar{y}^m - \hat{\beta}\bar{y}^d$, we can write
\begin{equation}
\hat{R}^2 = 1 - \frac{\sum_i(\bar{y}^m - \hat{\beta}\bar{y}^d + \hat{\beta}y_i^d - 
y_i^m)^2}{\sum_i (y_i^m - \bar{y}^m)^2}
\end{equation}
or 
\begin{equation}
\hat{R}^2 = 1 - \frac{\sum_i\left[\hat{\beta}(y_i^d  - \bar{y}^d) - (y_i^m - \bar{y}^m)
 \right]^2}{\sum_i (y_i^m - \bar{y}^m)^2}.
\end{equation}
In what follows, our objective is to understand how the above $\hat{R}^2$ relates to the ``true'' $R^2$ 
one would obtain from a regression of $y_i^m$ on $y_i^t$.

We first note that
\begin{equation}
\text{plim } \frac{N}{\sum_i (y_i^m - \bar{y}^m)^2} = \frac{1}{\sigma_m^2}
\label{eq:num}
\end{equation}
where $\sigma_m^2$ is the variance of the model-based poverty estimates.
Further, we can write
\begin{equation}
\text{plim }\frac{\sum_i\left[\hat{\beta}(y_i^d  - \bar{y}^d) - (y_i^m - \bar{y}^m)
 \right]^2}{N} =  \sigma_m^2 - \beta^2 \frac{\sigma_t^4}{\sigma_d^2}
 \label{eq:den}
\end{equation}
where $\sigma_t^2$ and $\sigma_d^2$ represent the variances of the true poverty rates and the direct 
estimates, respectively.
Combining Eqs.\ (\ref{eq:num}) and (\ref{eq:den}), we then have the following:
\begin{align}
\label{eq:r2}
\text{plim }\hat{R}^2 &= 1 - \frac{\sigma_m^2 - \beta^2 \frac{\sigma_t^4}{\sigma_d^2}}{\sigma_m^2} \\  \nonumber
& =  \beta^2 \frac{\sigma_t^4}{\sigma_d^2 \sigma_m^2} \\ \nonumber
& = \beta^2 \frac{\sigma_t^4}{\sigma_d^2 \sigma_m^2}\left(\frac{\sigma_t^2}{\sigma_t^2}\right) \\  \nonumber
& = \beta^2 \frac{\sigma_t^2}{ \sigma_m^2}\left(\frac{\sigma_t^2}{\sigma_d^2}\right) \\  \nonumber
& = R^2 \left(\frac{\sigma_t^2}{\sigma_d^2}\right)
\end{align}
where, for the last step in the above, we know that $R^2 = \beta^2 \sigma_t^2 / \sigma_m^2$, which is 
the ``true'' value of the $R^2$.

Given that $\sigma_d^2 = \sigma_t^2 + \sigma_\eta^2$, we can re-express Eq.\ (\ref{eq:r2}) as follows:
\begin{equation}
\hat{R}^2 = R^2 \left(\frac{\sigma_d^2 - \sigma_\eta^2}{\sigma_d^2}\right)
\end{equation}
where $\sigma_\eta^2$ denotes the variance of $\eta$.
In the above, we see that $\hat{R}^2 = R^2$ when $\sigma_\eta^2 = 0$ or, equivalently, when
$y_i^d = y_i^t$ for all $i$.
It is, however, well known that direct estimates provide noisy estimates of true poverty rates, 
meaning that $\sigma_\eta^2$ will generally be non-zero.
When $\sigma_\eta^2 > 0$, we then see that that $\hat{R}^2$ will be biased downward and the
magnitude of the bias will be determined by the size of $(\sigma_d^2 - \sigma_\eta^2)/ \sigma_d^2$.
We thus conclude that validation exercises based on regressions of the form in Eq.\ (\ref{eq:reg}) will
tend to mischaracterize the performance of model-based approaches to poverty mapping, with the 
resulting estimates of $R^2$ being systematically biased downward.





\end{document}